% Cote 119: offprints, reports, courses, CV and letters (typed and printed part): archive cote 119, site page 190 (PDF page 191).
% typescript. Draft compared: none.
% AI-transcribed by Claude Opus 5.5 (claude-opus-5-5, adaptive thinking, effort high), Message Batches API, batch r119_full (msgbatch_01DwZNm5Y7YsKBWyAN4sFs8w), 2026-11-13; not yet checked by eye.
\scanpage{119/190}
\mspage{- 9 -}

9.2. \uline{Points singuliers}.

Soit $f(z)$ une fonction holomorphe dans le disque pointé
$0 < |z| < \rho$ . On peut appliquer le théorème 8.2.2. pour obtenir un développement
de $f$ en série de Laurent. Si on désigne par $M(r)$ la borne supérieure de
$|f(z)|$ pour $|z| = r$ on obtient à partir de 8.2.3. les inégalités de Cauchy.
\[ (9.2.1) \qquad |a_n| \leq \frac{M(r)}{r^n} \quad n \in \mathbb{Z} \]

Pour que $f$ puisse être prolongée en une fonction holomorphe sur le
disque $|z| < \rho$ tout entier il faut évidemment que $f(z)$ soit bornée au voisi-%
nage de $0$. Cette condition est suffisante car en faisant tendre $r$ vers $0$
on voit que pour $n < 0$ on a nécessairement $a_n = 0$. Le développement de $f$ se
réduit donc à un développement en série entière.

Si la fonction $f$ ne peut être ainsi prolongée, on dit que $0$ est
un point singulier. Deux cas peuvent se présenter

- \uline{1° Cas} : Il n'existe qu'un nombre fini de $n < 0$ tels que $a_n \neq 0$
$f(z)$ est alors méromorphe. On dit que $0$ est un pôle.

- \uline{2° Cas} : Il existe une infinité de $n < 0$ tels que $a_n \neq 0$. Alors
$f(z)$ n'est pas méromorphe au voisinage de l'origine et on dit que $0$ est un
point singulier essentiel.

9.3. \uline{Théorème}. (Calcul des résidus).

Soit $U$ un ouvert connexe du plan complexe et soit $f$ une fonction
holomorphe dans $U$ sauf en des points isolés $z_k$ qui sont singuliers pour $f$.
$f$ est développable en série de Laurent en $(z-z_k)$ au voisinage de $z_k$ et l'on
appelle résidu de $f$ au point $z_k$ le coefficient de $\dfrac{1}{z - z_k}$ de ce développement.
Soit $\gamma$ un chemin fermé de $U$ ne passant par aucun point singulier et homotope
à un point dans $U$ . Alors il existe un nombre fini de points singuliers $z_k$
tels que $I(\gamma,z_k) \neq 0$ et l'on a :
\[ \int_\gamma f(z)\ dz = 2\pi i \ \sum_k \ I(\gamma,z_k)\ \mathrm{Res}(f,z_k) \]

\margin{190}
